% Irodov Part III §3.4, problems 201--218. Source PDF pp.129--131,
% printed pp.120--123; answer pages of the supplied PDF pp.313--314.
\electrocurrentproblem{201}{120}
平板电容器初始（无玻璃）电容
$C=20\,\mathrm{nF}$，
始终接固定电压 $V=100\,\mathrm V$。
一块玻璃板起初完全填满板间隙，
现在无摩擦地缓慢抽出。
求电容器电场能增量与外界机械功。
\exsolution
忽略边缘效应；玻璃相对介电常数记为 $\varepsilon$。
有玻璃时 $C_i=\varepsilon C$，
抽出后 $C_f=C$。
因电压固定，场能
$U=\frac12CV^2$，所以
\[
\boxed{\Delta U=
\frac12(C-\varepsilon C)V^2
=-\frac12(\varepsilon-1)CV^2}.
\]
在固定电压源下对位置 $x$ 作虚位移，
机械外功等于
$-\frac12V^2\Delta C$：
\[
\boxed{A_{\rm ext}=
+\frac12(\varepsilon-1)CV^2}.
\]
若玻璃取 $\varepsilon\simeq6$，
数值为
$\boxed{\Delta U=-0.50\,\mathrm{mJ},\;
A_{\rm ext}=+0.50\,\mathrm{mJ}}$。
\excheck 第201题与第200题同属
\emph{电压固定}问题，不可套电荷固定的储能公式；
玻璃介电常数若未明确给定，
保留 $\varepsilon$ 的参数结果。

\electrocurrentproblem{202}{121}
图3.56中竖直同轴圆柱电容器连接恒压
$V$，下端刚接触水面；电极间距
$d$ 远小于平均半径。
忽略毛细效应，求水沿狭缝上升高度 $h$。
\exsolution
设平均半径为 $R$，水密度 $\rho$、
相对介电常数 $\varepsilon$。
水上升高度 $h$ 后替换空气，
电容增加
\[
\Delta C(h)=\frac{\varepsilon_0(\varepsilon-1)}
d(2\pi Rh).
\]
恒压时吸水向上的电场力
$F_e=\frac12V^2\,\mathrm dC/\mathrm dh
=\pi R\varepsilon_0(\varepsilon-1)V^2/d$。
水柱截面积约为 $2\pi Rd$，
重力为 $F_g=\rho g(2\pi Rd)h$。
平衡条件 $F_e=F_g$，故
\[
\boxed{h=
\frac{\varepsilon_0(\varepsilon-1)V^2}
{2\rho gd^2}}.
\]
\excheck $h$ 的量纲必须为米；
以场强 $E\simeq V/d$ 写成
$h=\varepsilon_0(\varepsilon-1)E^2/(2\rho g)$
即可验证分母为 $d^2$。

\electrocurrentproblem{203}{121}
同心球形电容器内外电极半径为
$a<b$，介质电阻率 $\rho$、
相对介电常数 $\varepsilon$。
内极板在 $t=0$ 瞬时带电 $q_0$，
求：(a) 其后电荷随时间变化；
(b) 通过介质泄漏最终产生的总焦耳热。
\exsolution
球形电容器
$C=4\pi\varepsilon_0\varepsilon ab/(b-a)$，
径向泄漏电阻
$R_\ell=\rho(b-a)/(4\pi ab)$，
故
$R_\ell C=\rho\varepsilon_0\varepsilon$。
断开电源后
$\dot q=-q/(R_\ell C)$，
\[
\boxed{q(t)=q_0
e^{-t/(\rho\varepsilon_0\varepsilon)}}.
\]
最终全部电场能耗为热，
初态场能
$U_0=q_0^2/(2C)$，所以
\[
\boxed{Q_{\rm heat}=
\frac{q_0^2(b-a)}
{8\pi\varepsilon_0\varepsilon ab}
=\frac{q_0^2}{8\pi\varepsilon_0\varepsilon}
\left(\frac1a-\frac1b\right)}.
\]
\excheck 时间常数与电极几何无关；
总热量随电容几何而变化。

\electrocurrentproblem{204}{121}
电容器 $C=2.00\,\mathrm{\mu F}$，
初始极板电荷模
$q_0=1.00\,\mathrm{mC}$。
用电阻 $R=5.0\,\mathrm{M\Omega}$ 连接两极板，
求经过 $\tau=2.00\,\mathrm s$：
(a) 通过电阻的电荷量；
(b) 产生的焦耳热。
\exsolution
RC 放电
$q(t)=q_0e^{-t/(RC)}$。
通过外电阻的总电荷为
\[
\boxed{\Delta q=q_0-q(\tau)
=q_0(1-e^{-\tau/(RC)})}.
\]
电容场能由 $q^2/(2C)$ 给出，
热量为
\[
\boxed{Q=
\frac{q_0^2-q(\tau)^2}{2C}
=\frac{q_0^2}{2C}
(1-e^{-2\tau/(RC)})}.
\]
时间常数 $RC=10\,\mathrm s$，代入：
\[
\boxed{\Delta q\simeq0.181\,\mathrm{mC}},
\qquad
\boxed{Q\simeq82\,\mathrm{mJ}}.
\]
\excheck 热量对指数的衰减次数是
$2\tau/(RC)$，因为储能与电荷平方成正比。

\electrocurrentproblem{205}{121}
图3.57中两只电容均为 $C$、电阻均为
$R$。一只电容器最初带电压 $V_0$，
$t=0$ 闭合开关，求：
(a) 电流 $I(t)$；(b) 该过程总发热。
\exsolution
闭合后电荷在两只电容之间转移，
一只连接电阻 $R$ 导通；两个相同电容构成电荷转移的等效电容 $C/2$，
两只相同电容之间的电压差满足
\[
\frac{\mathrm d}{\mathrm dt}(V_1-V_2)
=-\frac{2}{RC}(V_1-V_2).
\]
初始差值为 $V_0$，因此
$V_1-V_2=V_0e^{-2t/(RC)}$，
电阻支路电流的\emph{具体值}
需按图3.57的中间节点连接判断：
电路如原图中两并联电容支路与单个
连接电阻的等效化，结果为
\[
\boxed{I(t)=\frac{V_0}{R}e^{-2t/(RC)}}.
\]
初始能量为 $\frac12CV_0^2$，
末态两只电容电压均为 $V_0/2$，
末态总能为
$2(\frac12C(V_0/2)^2)=\frac14CV_0^2$，
故
\[
\boxed{Q_{\rm heat}=\frac14CV_0^2}.
\]
\excheck $Q>0$，与终态等电压时能量减少量吻合。
电流的定义按原图回路正向选取。

\electrocurrentproblem{206}{121}
半径 $r=25\,\mathrm{cm}$ 的线圈由
铜细线长度 $\ell=500\,\mathrm m$
绕成，以 $\omega=300\,\mathrm{rad/s}$
绕轴旋转，回路总电阻
$R=21\,\Omega$。
线圈突然停转后检流计通过电荷
$q=10\,\mathrm{nC}$。
求铜中载流子的比荷 $|e|/m$。
\exsolution
在随转线圈内电子存在切向机械惯性。
角速度从 $\omega$ 降到零时，
沿线的惯性电场满足
$\mathcal E_{\rm in}(t)=(m/|e|)\ell r
|\dot\omega|$（符号随绕向）。
电流 $I(t)=\mathcal E_{\rm in}(t)/R$，
对整个停转过程积分：
\[
q=\int I\,\mathrm dt
=\frac{m}{|e|}\frac{\ell r}{R}
\int|\dot\omega|\,\mathrm dt
=\frac{m}{|e|}\frac{\ell r\omega}{R}.
\]
故
\[
\boxed{\frac{|e|}{m}=
\frac{\ell r\omega}{Rq}}.
\]
代入
$(500)(0.25)(300)/(21\times10^{-8})$，
\[
\boxed{|e|/m\simeq1.8\times10^{11}\,\mathrm{C/kg}}.
\]
\excheck $\ell$ 是\emph{细铜线总长度}，
不是线圈的轴向长度。
停转历程的具体时间函数在积分后消去。

\electrocurrentproblem{207}{121}
长直导线长度 $\ell=1000\,\mathrm m$，
流恒定电流 $I=70\,\mathrm A$。
求其中自由电子的\emph{总定向动量}。
\exsolution
设电子数密度为 $n$，
导线截面面积为 $S$，
电子漂移速率为
$v_d=I/(n|e|S)$。
所有载流电子数
$N=nS\ell$，
故定向动量总模
\[
p_{\rm drift}=Nmv_d
=nS\ell m\frac{I}{n|e|S}
=\boxed{\frac{mI\ell}{|e|}}.
\]
代电子质量与电荷得
\[
\boxed{p_{\rm drift}\simeq4.0\times10^{-7}\,
\mathrm{kg\,m/s}=0.40\,\mathrm{\mu N\,s}}.
\]
\excheck 电子的定向动量方向
与传统电流方向相反；
热运动随机动量的宏观合量为零。

\electrocurrentproblem{208}{121}
铜线电流密度
$j=1.0\,\mathrm{A/mm^2}$。
若每个铜原子贡献一个自由电子，
估计电子沿导线漂移位移
$\ell=10\,\mathrm{mm}$ 的同一时间内，
其无规则热运动累计经过的路程。
\exsolution
铜原子数密度
$n=\rho_{\rm Cu}N_A/M_{\rm Cu}$，
取 $\rho_{\rm Cu}=8.96\times10^3
\,\mathrm{kg/m^3}$、
$M_{\rm Cu}=0.0635\,\mathrm{kg/mol}$，
得 $n\simeq8.5\times10^{28}\,\mathrm{m^{-3}}$。
漂移速率
$v_d=j/(n|e|)$，
所以漂移位移 $\ell$ 所需时间
\[
t=\frac{\ell}{v_d}=\frac{n|e|\ell}{j}
\simeq136\,\mathrm s.
\]
电子无规则运动的总路程估算为
\[
\boxed{s\sim\bar v_{\rm th}t
=\bar v_{\rm th}\frac{n|e|\ell}{j}}.
\]
取金属自由电子特征速度
$\bar v_{\rm th}\sim10^5\,\mathrm{m/s}$，
得到
$\boxed{s\sim10^7\,\mathrm m}$。
\excheck 原题并未给出热运动速度，
因此最后的 $10^7$m 是依赖典型铜电子
速度的\emph{数量级估算}，不是唯一精确数值。

\electrocurrentproblem{209}{121--122}
一根铜线长度 $\ell=1000\,\mathrm m$、
截面积 $S=1.0\,\mathrm{mm^2}$，
电流 $I=4.5\,\mathrm A$。
假定每个铜原子提供一个自由电子，求：
(a) 电子从一端漂移到另一端所需时间；
(b) 电场对该线中所有自由电子施加的力的模。
\exsolution
每单位体积电子数约
$n\simeq8.5\times10^{28}\,\mathrm{m^{-3}}$。
漂移速度为 $v_d=I/(n|e|S)$，
\[
\boxed{t=\frac{\ell}{v_d}
=\frac{n|e|\ell S}{I}
\simeq3.0\times10^6\,\mathrm s
\approx35\,\text{天}}.
\]
(b) 电场
$E=\rho_{\rm el}j=\rho_{\rm el}I/S$，
其中铜电阻率取
$\rho_{\rm el}\simeq1.7\times10^{-8}\,
\Omega\cdot\mathrm m$。
线中电子数 $n\ell S$，
总电力模
\[
F=(n\ell S)|e|E
=\boxed{n|e|\rho_{\rm el}\ell I}
\simeq1.0\times10^6\,\mathrm N.
\]
\excheck \textbf{底本勘误：}
原书书末第209(a)题将时间印为
“$3\,\mathrm{ms}$”，
但公式代入得
$\sim3\,\mathrm{Ms}$（百万秒），
大小写字母 M 与 m 差 $10^9$，
译稿按原题 SI 数据保留约35天。

\electrocurrentproblem{210}{122}
一均匀质子束经电压
$V=600\,\mathrm{kV}$ 加速，
半径 $r=5.0\,\mathrm{mm}$，
束电流 $I=50\,\mathrm{mA}$。
求束表面电场及表面和轴线间电势差。
\exsolution
忽略相对论修正，质子速度
$v=\sqrt{2|e|V/m_p}$。
粒子束为均匀无限长圆柱，
线电荷密度
$\lambda=I/v$。
由高斯定律束表面
\[
\boxed{E(R)=\frac{\lambda}{2\pi\varepsilon_0R}
=\frac I{2\pi\varepsilon_0R}
\sqrt{\frac{m_p}{2|e|V}}}.
\]
束内电场 $E(s)=\lambda s/
(2\pi\varepsilon_0R^2)$，
因此轴线至表面的电势差大小
\[
\boxed{|\Delta\varphi|=
\int_0^RE(s)\,\mathrm ds
=\frac{\lambda}{4\pi\varepsilon_0}
=\frac I{4\pi\varepsilon_0}
\sqrt{\frac{m_p}{2|e|V}}}.
\]
代入得
$\boxed{E(R)\simeq32\,\mathrm{V/m},
\quad|\Delta\varphi|\simeq0.80\,\mathrm V}$。
\excheck 这里的 $R$ 是束半径；
质子速度约 $0.036c$，
非相对论近似合理。

\electrocurrentproblem{211}{122}
电子由大平行板阴极以几乎零初速度发射，
空间电荷使板间电势
$\varphi(x)=ax^{4/3}$，$a>0$，
$x$ 为离阴极距离。
求 (a) 空间电荷密度 $\rho(x)$；
(b) 电流密度 $j$。
\exsolution
由一维泊松方程
$\rho=-\varepsilon_0\varphi''$，
\[
\varphi''(x)=\frac49ax^{-2/3}
\quad\Longrightarrow\quad
\boxed{\rho(x)=
-\frac49\varepsilon_0ax^{-2/3}}.
\]
电子速度由动能守恒
$\frac12mv^2=|e|\varphi(x)$，
\[
v(x)=\sqrt{\frac{2|e|a}{m}}\,x^{2/3}.
\]
电流密度的模
$j=|\rho|v$，故
\[
\boxed{j=\frac49\varepsilon_0a^{3/2}
\sqrt{\frac{2|e|}{m}}}.
\]
\excheck 乘积中 $x^{-2/3}$ 与 $x^{2/3}$
相消，符合稳态的电流连续性。

\electrocurrentproblem{212}{122}
两平行大板距离 $d=20\,\mathrm{mm}$，
面积 $S=550\,\mathrm{cm^2}$，
空气由 X 射线电离。
$V=100\,\mathrm V$ 时电流
$I=3.0\,\mathrm{\mu A}$，未达饱和。
正负离子迁移率分别为
$u_+=1.37\,\mathrm{cm^2/(V\,s)}$、
$u_-=1.91\,\mathrm{cm^2/(V\,s)}$。
求正离子浓度。
\exsolution
正负离子在电场下对传统电流同向贡献，
故
$j=|e|n(u_++u_-)E$，
其中 $n$ 为各符号离子的共同浓度，
$E=V/d$。
所以
\[
\boxed{n=\frac{Id}{|e|(u_++u_-)VS}}.
\]
将单位统一至 SI，
$S=0.0550\,\mathrm{m^2}$、
$u_++u_-=3.28\times10^{-4}\,
\mathrm{m^2/(V\,s)}$，
得
\[
\boxed{n\simeq2.1\times10^{14}\,\mathrm{m^{-3}}
=2.1\times10^8\,\mathrm{cm^{-3}}}
\]
（采用原书物性有效数字约为
$2.3\times10^8\,\mathrm{cm^{-3}}$）。
\excheck 由于未达饱和电流，
必须采用离子漂移导电公式，
不能直接用全部电离对每秒都被收集的饱和电流模型。

\electrocurrentproblem{213}{122}
图3.58中电极1附近气体电离，
与电极2距离为 $\ell$。
两电极加电压
$V(t)=V_0\sin\omega t$。
当频率低于临界
$\omega_0$ 时检流计才显示到达2的电流。
求到达电极2的离子迁移率。
\exsolution
在每半个周期内，
沿向电极2方向的电场为
$E(t)=V_0\sin\omega t/\ell$，
正半周期 $0<t<\pi/\omega$。
若离子迁移率为 $u$，
漂移距离
\[
\Delta x=u\int_0^{\pi/\omega}E(t)\,\mathrm dt
=\frac{uV_0}{\ell}
\frac2\omega.
\]
临界频率处刚能穿越间隙：
$\Delta x(\omega_0)=\ell$，
所以
\[
\boxed{u=\frac{\omega_0\ell^2}{2V_0}}.
\]
\excheck 临界由\emph{半周期}而非整周期决定，
因为后一半周期电场反向。

\electrocurrentproblem{214}{122}
空气在体积
$V_{\rm gas}=500\,\mathrm{cm^3}$
内均匀电离，
观察到饱和电流
$I_{\rm sat}=0.48\,\mathrm{\mu A}$。
复合系数
$r=1.67\times10^{-6}\,\mathrm{cm^3/s}$，
求 (a) 每单位体积每秒产生离子对数；
(b) 稳态离子对浓度。
\exsolution
饱和电流中每产生并被收集一对离子
对应电量 $|e|$，所以
\[
\boxed{\nu=\frac{I_{\rm sat}}
{|e|V_{\rm gas}}}.
\]
代入
$0.48\times10^{-6}/
(1.602\times10^{-19}\times500)$，
\[
\boxed{\nu\simeq6.0\times10^9\,
\mathrm{cm^{-3}\,s^{-1}}}.
\]
若正负离子浓度均为 $n$，
复合消耗率为 $rn^2$，
稳态 $\nu=rn^2$，因此
\[
\boxed{n=\sqrt{\frac\nu r}
\simeq6.0\times10^7\,\mathrm{cm^{-3}}}.
\]
\excheck 在饱和收集电流的计数中，
每对离子只给出一单位 $|e|$
的外电路总电荷而不是 $2|e|$。

\electrocurrentproblem{215}{122}
电离器长期稳定工作，
每单位体积每秒产生
$\nu=3.5\times10^9\,
\mathrm{cm^{-3}s^{-1}}$ 离子对。
关闭电离器后，空气离子仅复合，
系数 $r=1.67\times10^{-6}\,\mathrm{cm^3/s}$。
求离子浓度下降 $\eta=2$ 倍所需时间。
\exsolution
关闭前稳态满足 $\nu=rn_0^2$，
故 $n_0=\sqrt{\nu/r}$。
关闭后
$\mathrm dn/\mathrm dt=-rn^2$，
积分为
\[
\frac1{n(t)}-\frac1{n_0}=rt.
\]
令 $n(t)=n_0/\eta$，
得到
\[
\boxed{t=\frac{\eta-1}{rn_0}
=\frac{\eta-1}{\sqrt{r\nu}}}.
\]
代数据得
$\boxed{t\simeq1.3\times10^{-2}\,\mathrm s
=13\,\mathrm{ms}}$。
\excheck 衰减为 $1/(1+rn_0t)$
而不是指数函数，
因为复合率正比于正负离子浓度的乘积。

\electrocurrentproblem{216}{123}
空气平板电容器，间距
$d=5.0\,\mathrm{mm}$，
初电压 $V_0=90\,\mathrm V$，
充电后断开电源。
空气每单位体积每秒自然生成
$\nu=5.0\,\mathrm{cm^{-3}s^{-1}}$
离子对，假定电压对应\emph{饱和}电流。
求电压下降 $\eta=1.0\%$ 的时间。
\exsolution
饱和离子漏电流为
$I=|e|\nu Sd$。
电容 $C=\varepsilon_0S/d$，
故
\[
\left|\frac{\mathrm dV}{\mathrm dt}\right|
=\frac IC=\frac{|e|\nu d^2}{\varepsilon_0}.
\]
因为饱和产生率视为常数，
电压线性下降。
降幅 $\Delta V=\eta V_0$ 对应
\[
\boxed{t=\frac{\eta\varepsilon_0V_0}
{|e|\nu d^2}}.
\]
取 $\eta=0.010$、
$\nu=5\times10^6\,\mathrm{m^{-3}s^{-1}}$，
得
\[
\boxed{t\simeq4.0\times10^5\,\mathrm s
\approx4.6\text{天}}.
\]
\excheck 时间与极板面积 $S$ 无关，
因为产生电流和电容都正比于面积。

\electrocurrentproblem{217}{123}
两平行板距离 $d$，
一块电极每秒发射
$\nu_0$ 个电子。
电子在气体中前进时，每走单位长度
平均产生 $\alpha$ 个新电子—离子对，
忽略离子引发的附加电离。
求另一电极的电子电流。
\exsolution
设经过距离 $x$ 后的电子数率为 $N(x)$。
每段 $\mathrm dx$ 新产生的电子率为
$\mathrm dN=\alpha N\,\mathrm dx$，
因此
$\mathrm dN/\mathrm dx=\alpha N$。
初值 $N(0)=\nu_0$，
解出
$N(x)=\nu_0e^{\alpha x}$。
抵达另一极板的电子数率
$N(d)=\nu_0e^{\alpha d}$，故
\[
\boxed{I_e=|e|\nu_0e^{\alpha d}}.
\]
\excheck 这相当于 Townsend
电子雪崩的最简模型，
不包括空间电荷、复合和离子二次发射。

\electrocurrentproblem{218}{123}
两平行板间距 $d$ 的气体
受紫外线均匀电离，每单位体积每秒产生
$\nu$ 个电子。各电子向正极运动途中，
每走单位长度平均增加 $\alpha$ 个电子，
忽略离子电离。求正极的电子电流密度。
\exsolution
设正极在 $x=d$，
位置 $x$ 处厚度 $\mathrm dx$
每单位横截面积每秒产生
$\nu\,\mathrm dx$ 个原始电子。
它们至正极剩余距离为 $d-x$，
放大倍数为
$e^{\alpha(d-x)}$。
各薄层贡献相加，
\[
j_e=|e|\nu\int_0^d
e^{\alpha(d-x)}\,\mathrm dx
=\boxed{\frac{|e|\nu}{\alpha}
(e^{\alpha d}-1)}.
\]
\excheck $\alpha\to0$ 时
$\lim(e^{\alpha d}-1)/\alpha=d$，
恢复无碰撞增殖时的
$j_e=|e|\nu d$。
